\chapter{词汇表}\label{s:gloss}

\begin{description}

\gitem{g:free-range-learner}{散养型学习者} 在体制内课堂之外学习、没有规定的作业和统一课程的人。（使用这个词的人有时会把课堂里的学习者称为「圈养型学习者」，但我们不这么叫，因为那太失礼了。）

\gitem{g:end-user-teacher}{终端用户教师} 仿照\grefcross{g:end-user-programmer}{终端用户程序员}的说法，指经常在教书、但教书并非其主业的人；他们几乎没有或完全没有教育学背景，并且可能在体制内课堂之外工作。

\gitem{g:impostor-syndrome}{冒充者综合征} 对自己的成就感到不安，表现为害怕被揭穿是个骗子。

\gitem{g:novice}{新手} 尚未对某一领域建立起可用心智模型的人。另见\grefcross{g:competent-practitioner}{胜任者}和\grefcross{g:expert}{专家}。

\gitem{g:competent-practitioner}{胜任者} 在正常情况下以正常的努力完成常规任务的人。另见\grefcross{g:novice}{新手}和\grefcross{g:expert}{专家}。

\gitem{g:expert}{专家} 能够诊断和处理不寻常情况、知道常规规则何时不适用、并且往往能直接认出解而不必一步步推演的人。另见\grefcross{g:competent-practitioner}{胜任者}和\grefcross{g:novice}{新手}。

\gitem{g:mental-model}{心智模型} 对某一问题领域的关键要素及关系的简化表示，其精确程度足以支撑解决问题。

\gitem{g:tutorial}{教程} 旨在帮助某人加深对某一学科总体理解的课程。

\gitem{g:manual}{手册} 旨在帮助已经理解某一学科的人补齐（或提醒自己）细节的参考资料。

\gitem{g:expertise-reversal}{专长逆转效应} 对新手有效的教学，对胜任者或专家反而变得无效的现象。

\gitem{g:formative-assessment}{形成性评价} 在课程进行过程中实施的评价，目的是让学习者和教师都得到关于真实理解程度的反馈。另见\grefcross{g:summative-assessment}{总结性评价}。

\gitem{g:summative-assessment}{总结性评价} 在一门课程结束时实施的评价，用来判断想达成的学习是否已经发生。

\gitem{g:diagnostic-power}{诊断力} 一个问题或练习的错误答案，能在多大程度上告诉教师某位学习者具体有哪些误解。

\gitem{g:plausible-distractor}{似真干扰项} 多项选择题中看起来像是正确答案的错误答案或不那么理想的答案。另见\grefcross{g:diagnostic-power}{诊断力}。

\gitem{g:active-teaching}{主动教学} 一种教学方式，教师在授课过程中会根据从学习者那里获得的新信息采取行动（例如动态地更换例子，或调整内容的原定顺序）。另见\grefcross{g:passive-teaching}{被动教学}。

\gitem{g:passive-teaching}{被动教学} 一种教学方式，教师在课上不调整进度或例子，也不对学习者的反馈作出其他反应。另见\grefcross{g:active-teaching}{主动教学}。

\gitem{g:concept-inventory}{概念清单} 一种用来判断学习者对某个领域理解程度的测验。与大多数由教师自编的测验不同，概念清单建立在大量研究和验证的基础上。

\gitem{g:computational-thinking}{计算思维} 以受编程启发的思路来思考如何解决问题（不过这个词还有许多别的用法）。

\gitem{g:notional-machine}{观念机器} 对某一家族的程序如何执行所做的一种通用、简化的模型。

\gitem{g:intuition}{直觉} 无需任何明显的自觉推理就能立刻理解某事的能力。

\gitem{g:fluid-representation}{流畅表征} 在不同的解题模型之间快速切换的能力。

\gitem{g:expert-blind-spot}{专家盲区} 专家难以对第一次接触某些概念或做法的初学者产生共情。

\gitem{g:concept-map}{概念图} 一种心智模型的图像，其中概念是图中的节点，概念之间的关系是有（标注的）弧。

\gitem{g:externalized-cognition}{外化认知} 借助图形、实物或言语辅助来增强思考。

\gitem{g:long-term-memory}{长期记忆} 记忆中长期存储信息的那部分。长期记忆容量极大，但速度慢。另见\grefcross{g:short-term-memory}{短期记忆}。

\gitem{g:persistent-memory}{持久记忆} 见\grefcross{g:long-term-memory}{长期记忆}。

\gitem{g:short-term-memory}{短期记忆} 记忆中可以由意识直接访问、短暂存储信息的那部分。

\gitem{g:working-memory}{工作记忆} 见\grefcross{g:short-term-memory}{短期记忆}。

\gitem{g:chunking}{组块化} 把相关概念归到一起，使它们能作为一个单元被存储和处理的做法。

\gitem{g:deliberate-practice}{刻意练习} 在做一项任务的同时观察自己的表现，以便提高能力的做法。

\gitem{g:reflective-practice}{反思性实践} 见\grefcross{g:deliberate-practice}{刻意练习}。

\gitem{g:split-attention-effect}{注意分散效应} 当学习者必须在同一信息的多种同时呈现之间分配注意力时（例如字幕和配音），学习效果会下降。

\gitem{g:inquiry-based-learning}{探究式学习} 允许学习者提出自己的问题、设定自己的目标、并自行找到学习某一学科的路径的做法。

\gitem{g:cognitive-load}{认知负荷} 解决问题所需的心理投入。认知负荷理论把它分为\grefcross{g:intrinsic-load}{内在}、\grefcross{g:germane-load}{关联}和\grefcross{g:extraneous-load}{外在}负荷，并认为当关联负荷和外在负荷降低时，人学得最快。

\gitem{g:intrinsic-load}{内在负荷} 吸收新信息所需的\grefcross{g:cognitive-load}{认知负荷}。

\gitem{g:germane-load}{关联负荷} 把新信息和旧信息联系起来所需的\grefcross{g:cognitive-load}{认知负荷}。

\gitem{g:extraneous-load}{外在负荷} 任何会分散学习注意力的\grefcross{g:cognitive-load}{认知负荷}。

\gitem{g:parsons-problem}{帕森斯问题} 由 Dale Parsons 等人开发的一种评估方法，学习者重新排列给定的材料，从而构造出问题的正确答案~\cite{Pars2006}。

\gitem{g:faded-example}{渐隐示例} 一系列例子，其中被留空的关键步骤逐步增多。另见\grefcross{g:scaffolding}{脚手架}。

\gitem{g:scaffolding}{脚手架} 为处于早期阶段的学习者提供的额外材料，帮助他们解决问题。

\gitem{g:cognitive-apprenticeship}{认知学徒制} 一种学习理论，强调师傅在具体情境中把技能和洞见传授给学徒的过程。

\gitem{g:subgoal-labeling}{子目标标注} 为问题解决过程的分步描述中的各个步骤命名。

\gitem{g:minimal-manual}{极简手册} 一种培训方法，把每项任务拆成一页纸的说明，并同时解释如何诊断和纠正常见错误。

\gitem{g:educational-psychology}{教育心理学} 研究人如何学习的学科。另见\grefcross{g:instructional-design}{教学设计}。

\gitem{g:instructional-design}{教学设计} 为特定受众设计并评价具体课程的技艺。另见\grefcross{g:educational-psychology}{教育心理学}。

\gitem{g:cognitivism}{认知主义} 一种学习理论，认为心理状态和过程可以而且必须被纳入学习模型之中。另见\grefcross{g:behaviorism}{行为主义}。

\gitem{g:situated-learning}{情境学习} 一种学习模型，关注人从新来者转变为\grefcross{g:community-of-practice}{实践共同体}中被接纳的成员这一过程。

\gitem{g:behaviorism}{行为主义} 一种学习理论，其核心原则是刺激与反应，目标是在不借助内部心理状态或其他不可观测因素的情况下解释行为。另见\grefcross{g:cognitivism}{认知主义}。

\gitem{g:constructivism}{建构主义} 一种学习理论，把学习者看作在主动地建构知识。

\gitem{g:connectivism}{联通主义} 一种学习理论，认为知识是分布式的，学习就是建立、生长和修剪连接的过程，并强调互联网使学习具有的社会性。

\gitem{g:passive-learning}{被动学习} 一种教学方式，学习者只是读、听或看，而不立刻运用新知识。被动学习的效果不如\grefcross{g:active-learning}{主动学习}。

\gitem{g:active-learning}{主动学习} 一种教学方式，学习者通过讨论、解决问题、案例分析以及其他需要他们在当下反思并运用新信息的活动来接触材料。另见\grefcross{g:passive-learning}{被动学习}。

\gitem{g:metacognition}{元认知} 对思考的思考。

\gitem{g:transfer-of-learning}{学习迁移} 把在某一情境中学到的知识应用到另一情境中的问题。另见\grefcross{g:near-transfer}{近迁移}和\grefcross{g:far-transfer}{远迁移}。

\gitem{g:near-transfer}{近迁移} 在密切相关领域之间的\grefcross{g:transfer-of-learning}{学习迁移}，例如因为做分数练习而加深了对小数的理解。

\gitem{g:far-transfer}{远迁移} 在相距甚远的领域之间的\grefcross{g:transfer-of-learning}{学习迁移}，例如因为下棋而提高了数学能力。

\gitem{g:transfer-appropriate-processing}{迁移适宜加工} 当练习所用的活动与测验所用的活动相似时，回忆效果会有所提高。

\gitem{g:read-cover-retrieve}{阅读—遮盖—回忆} 一种学习练习：学习者第一遍看材料时遮住关键事实或术语，第二遍时检查自己能回忆起多少。

\gitem{g:hypercorrection}{超校正效应} 一个人在测验中越是坚信自己的答案是对的，一旦发现其实错了，就越有可能不再犯同样的错。

\gitem{g:automaticity}{自动化} 在不需要专注于低层次细节的情况下完成一项任务的能力。

\gitem{g:flow}{心流} 完全沉浸在某一活动中的感觉；常与高效率联系在一起。

\gitem{g:contributing-student}{贡献型学生教学法} 让学习者制作作品，为其他人的学习作出贡献。

\gitem{g:calibrated-peer-review}{校准式同伴互评} 先让学习者把自己对样卷的评语与教师的评语作比较，之后才允许他们去评阅同伴的作业。

\gitem{g:guided-notes}{引导式笔记} 教师事先准备好的笔记，用来提示学习者在讲座或讨论中对关键信息作出回应。

\gitem{g:test-driven-development}{测试驱动开发} 一种软件开发实践，程序员先写测试，以便给自己设定具体目标，并厘清自己对「完成」意味着什么的理解。

\gitem{g:backward-design}{逆向设计} 一种教学设计方法，它从总结性评价往回推，推到期中形成性评价，再推到课程内容。

\gitem{g:teaching-to-the-test}{应试教学} 任何以让学生通过标准化考试为目的、而非为了真正学习的「教育」方法。

\gitem{g:learner-persona}{学习者画像} 对某一课程典型目标学习者的简要描述，包括其大致背景、已经会什么、想做什么、这门课将如何帮到他们，以及他们可能有的特殊需求。

\gitem{g:learning-objective}{学习目标} 一门课程想要达成的结果。

\gitem{g:learning-outcome}{学习成效} 一门课程实际达成的结果。

\gitem{g:blooms-taxonomy}{布鲁姆分类法} 一种将理解分为六个层级的分类法，已被广泛采用，其层级为\emph{知识}、\emph{理解}、\emph{应用}、\emph{分析}、\emph{综合}和\emph{评价}。另见\grefcross{g:finks-taxonomy}{芬克分类法}。

\gitem{g:finks-taxonomy}{芬克分类法} 一种将理解分为六个非层级类别的分类法，最初由~\cite{Fink2013}提出，其类别为\emph{基础知识}、\emph{应用}、\emph{整合}、\emph{人的维度}、\emph{关怀}和\emph{学会学习}。另见\grefcross{g:blooms-taxonomy}{布鲁姆分类法}。

\gitem{g:reusability-paradox}{可复用性悖论} 认为课程的某部分越可复用，其教学效果就越差。

\gitem{g:zpd}{最近发展区（ZPD）} 包含人们目前还不能独立解决、但在更有经验的导师帮助下可以解决的一类问题。另见\grefcross{g:productive-failure}{有效失败}。

\gitem{g:content-knowledge}{内容知识} 一个人对某一学科的理解。另见\grefcross{g:general-pedagogical-knowledge}{一般教学法知识}和\grefcross{g:pedagogical-content-knowledge}{学科教学知识}。

\gitem{g:general-pedagogical-knowledge}{一般教学法知识} 一个人对教学的一般原理的理解。另见\grefcross{g:content-knowledge}{内容知识}和\grefcross{g:pedagogical-content-knowledge}{学科教学知识}。

\gitem{g:pedagogical-content-knowledge}{学科教学知识（PCK）} 关于如何教授某一特定学科的理解，即按什么顺序引入各个主题最好、该用哪些例子。另见\grefcross{g:content-knowledge}{内容知识}和\grefcross{g:general-pedagogical-knowledge}{一般教学法知识}。

\gitem{g:end-user-programmer}{终端用户程序员} 并不认为自己是程序员，却仍然在编写和调试软件的人，例如为绘图工具编写复杂宏的艺术家。

\gitem{g:cs1}{CS1} 大学计算机科学入门课程，通常持续一个学期，重点讲变量、循环、函数以及其他基础机制。

\gitem{g:cs2}{CS2} 大学计算机科学的第二门课程，通常介绍栈、队列和字典等基本数据结构。

\gitem{g:cs0}{CS0} 面向非专业学生、几乎没有或完全没有编程先修经验的大学入门计算课程。

\gitem{g:objects-first}{对象优先} 一种教授编程的方法，很早就介绍对象和类。

\gitem{g:cs-unplugged}{不插电计算机科学} 一种用非编程的例子和实物来介绍计算概念的教学风格。

\gitem{g:live-coding}{现场编程} 随着课程推进，当着学习者的面编写软件来教授编程。

\gitem{g:twitch-coding}{直播编程} 由一群人一起来决定下一步该往程序里加什么，可能是逐刻决定，也可能是逐行决定。

\gitem{g:direct-instruction}{直接教学} 一种围绕精细设计的课程、通过规定脚本进行讲授的教学方法。

\gitem{g:jugyokenkyu}{授业研究} 字面意思是「课例研究」，是一套做法，包括让教师定期互相观摩、讨论课程，以分享知识、改进技能。

\gitem{g:demonstration-lesson}{示范课} 一种安排好的课，由一位教师对真实学习者授课，其他教师从旁观摩，以学习新的教学技巧。

\gitem{g:peer-instruction}{同伴教学} 一种教学方法，教师提出一个问题，学习者先给出自己的第一个答案，再与同伴讨论答案，然后（修正后）给出一个答案。

\gitem{g:co-teaching}{协同教学} 与另一位教师一起在同一间教室里教学。

\gitem{g:dunning-kruger-effect}{邓宁—克鲁格效应} 对某一学科只懂一点点的人，倾向于错误地高估自己对该学科的理解。

\gitem{g:false-beginner}{假性新手} 以前学过一门语言、现在又重新学的人。假性新手的起点和真正的新手一样（即前测会显示同样的水平），但进步速度可以快得多。

\gitem{g:absolute-beginner}{绝对新手} 从没接触过相关概念或材料的人。这个词是用来和\grefcross{g:false-beginner}{假性新手}相区分的。

\gitem{g:preparatory-privilege}{预备性特权} 来自一个能为某项学习任务提供比他人更多准备的背景所带来的优势。

\gitem{g:pair-programming}{结对编程} 一种软件开发实践，两位程序员共用一台电脑。一位程序员（驾驶员）负责打字，另一位（领航员）实时提出意见和想法。结对编程常被用作编程课上的教学实践。

\gitem{g:minute-cards}{一分钟卡片} 一种反馈技巧，学习者花一分钟写下关于一节课的一件正面的事（例如学到的一点）和一件负面的事（例如一个仍未得到回答的问题）。

\gitem{g:think-pair-share}{思考—结对—分享} 一种协作方法：每个人先独自思考一个问题或难题，然后与一位同伴结对交流想法，最后每组派一人向全班汇报。

\gitem{g:extrinsic-motivation}{外在动机} 由报酬或害怕惩罚等外部奖赏所驱动。另见\grefcross{g:intrinsic-motivation}{内在动机}。

\gitem{g:intrinsic-motivation}{内在动机} 由享受某项任务本身或为做而做的满足感所驱动。另见\grefcross{g:extrinsic-motivation}{外在动机}。

\gitem{g:authentic-task}{真实任务} 包含学习者在真实（非课堂）情境中会做之事的要素的任务。要算得上真实，任务应当要求学习者自己构造答案，而不是在给定的选项之间挑选，并且使用他们在现实生活中会用的同样工具和数据。

\gitem{g:tangible-artifact}{有形造物} 学习者可以动手去做、并且其状态能反馈学习者进展、帮助其诊断错误的东西。

\gitem{g:learned-helplessness}{习得性无助} 人反复遭受无法逃避的负面反馈后，即使有机会逃离也不再尝试逃离的情形。

\gitem{g:productive-failure}{有效失败} 一种情形：学习者被刻意给到用他们现有知识无法解决的问题，必须走出去获取新信息才能取得进展。另见\grefcross{g:zpd}{最近发展区}。

\gitem{g:fixed-mindset}{固定型思维} 认为能力是天生的、失败是因为缺少某种必备属性。另见\grefcross{g:growth-mindset}{成长型思维}。

\gitem{g:growth-mindset}{成长型思维} 认为能力来自练习。另见\grefcross{g:fixed-mindset}{固定型思维}。

\gitem{g:stereotype-threat}{刻板印象威胁} 一种情形：人们感到自己有被按照其社会群体的刻板印象来评判的风险。

\gitem{g:inclusivity}{包容性} 积极努力地接纳背景和需求各异的人。

\gitem{g:community-representation}{社群表征} 在学习活动中运用文化资本，凸显学习者的社会身份、历史和社群网络。

\gitem{g:computational-integration}{计算性整合} 用计算来重新实现既有的文化造物，例如用计算机绘图工具创作传统图案的变体。

\gitem{g:deficit-model}{缺陷模型} 认为某些群体在计算机（或其他某个领域）中代表性不足，是因为其成员缺少某种属性或品质的观点。

\gitem{g:mooc}{大规模开放在线课程（MOOC）} 为大规模选课和异步学习而设计的在线课程，通常使用录制视频和自动评分。

\gitem{g:personalized-learning}{个性化学习} 自动调整课程，以满足每位学习者的需求。

\gitem{g:lms}{学习管理系统（LMS）} 用于跟踪选课、作业提交、成绩以及其他正规课堂教学事务性环节的应用程序。

\gitem{g:flipped-classroom}{翻转课堂} 学习者利用自己的时间观看录制的课程，而课堂时间用来做题和答疑。

\gitem{g:fuzz-testing}{\emph{fuzz testing}} 一种基于生成并输入随机数据的软件测试技术。

\gitem{g:hashing}{哈希} 从数据生成一个浓缩的伪随机数字密钥；任何特定输入都会产生相同的输出，而不同输入极有可能产生不同的输出。

\gitem{g:legitimate-peripheral-participation}{合法的边缘性参与} 新来者在简单、低风险任务中的参与，这种参与被\grefcross{g:community-of-practice}{实践共同体}认可为有效的贡献。

\gitem{g:community-of-practice}{实践共同体} 一个自我延续的群体，成员共享并共同发展某一门手艺，例如编织爱好者、音乐人或程序员。另见\grefcross{g:legitimate-peripheral-participation}{合法的边缘性参与}。

\gitem{g:commons}{共有资源} 由某个社群按照他们自己演化并采纳的规则共同管理的东西。

\gitem{g:service-board}{服务型委员会} 其成员在组织中承担实际工作角色的委员会。

\gitem{g:governance-board}{治理委员会} 主要职责是聘用、监督并在必要时解雇负责人的委员会。

\gitem{g:ego-depletion}{自我损耗} 因长时间或高强度使用而导致的自我控制能力减弱。近期研究未能证实其存在。

\gitem{g:marketing}{市场营销} 从别人的视角看事情、理解他们的需求和想要什么，并设法满足这些需求的技艺。

\gitem{g:elevator-pitch}{电梯演讲} 对某个想法、项目、产品或人的简短描述，能在短短几秒内讲完并让人听懂。

\gitem{g:brand}{品牌} 人们对某个产品名称或身份所产生的联想。

\gitem{g:positioning}{定位} 把某个品牌与其他类似的品牌区分开来的东西。

\gitem{g:seo}{搜索引擎优化} 通过让网页更容易被搜索引擎找到、或显得对搜索引擎更重要，来提高网站流量的数量和质量。

\gitem{g:implementation-science}{实施科学} 研究如何把研究成果转化为日常临床实践的学科。

\gitem{g:conversational-programmer}{对话型程序员} 需要懂一些计算知识以便与程序员进行有意义的交流，但自己并不打算编程的人。

\gitem{g:pull-request}{拉取请求（pull request）} 对某个 GitHub 仓库提出的一整套改动建议，可以被审阅、更新，并最终合并。

\end{description}
